\documentclass[a4paper, serif, 10pt]{moderncv}

%% moderncv
% color and style
\moderncvcolor{black}
\moderncvstyle{banking}

%% page layout/margins
\usepackage[top=2.5cm, bottom=2.5cm, left=1.5cm, right=1.5cm]{geometry}

\nopagenumbers{}

%% Babel config for German language
% ref:
% - https://latex3.github.io/babel/guides/locale-german.html
\usepackage[ngerman]{babel}

%% fontspec
\usepackage{fontspec}

\IfFontExistsTF{Linux Libertine}{
  \setmainfont[Ligatures={TeX, Common}, Numbers=OldStyle]{Linux Libertine}
}{}
\IfFontExistsTF{Linux Libertine O}{
  \setmainfont[Ligatures={TeX, Common}, Numbers=OldStyle]{Linux Libertine O}
}{}

%% CTeX
% CJK support, used to show CJK characters in the resume
%
% - heading=false: disable ctex heading
% - punct=kaiming: use kaiming punctuations styles for CJK
% - scheme=plain: use plain scheme, do not override \normalsize font size
% - space=auto: space settings for CJK characters
%
% CJK fonts are resolved per locale language, e.g. Noto Serif/Sans CJK JP
% for Japanese resumes, with fallback to the next region if not installed
%
% ref:
% - http://ctan.mirrorcatalogs.com/language/chinese/ctex/ctex.pdf
\usepackage[UTF8, heading=false, punct=kaiming, scheme=plain, space=auto]{ctex}

\IfFontExistsTF{Noto Serif CJK SC}{
  \setCJKmainfont{Noto Serif CJK SC}
}{}
\IfFontExistsTF{Noto Sans CJK SC}{
  \setCJKsansfont{Noto Sans CJK SC}
}{}

%% line spacing
\usepackage{setspace}
\setstretch{1.125}

%% URL styling - use normal text instead of monospace
\urlstyle{same}

%% Patch \httplink and \httpslink to handle full URLs
% The original moderncv macros always prepend http:// or https:// to URLs.
% This causes issues when the URL already has a protocol (e.g., https://example.com)
% resulting in malformed URLs like https://https://example.com.
% This patch checks if the URL already contains "://" and uses it directly if so.
% Note: We use \str_set:Nx (x-expansion) to expand macros like \@homepage first.
\ExplSyntaxOn
\str_new:N \g_yamlresume_url_str

\RenewDocumentCommand{\httpslink}{O{}m}{
  \str_set:Nx \g_yamlresume_url_str {#2}
  \str_if_in:NnTF \g_yamlresume_url_str {://}
    { \url{#2} }
    { \url{https://#2} }
}

\RenewDocumentCommand{\httplink}{O{}m}{
  \str_set:Nx \g_yamlresume_url_str {#2}
  \str_if_in:NnTF \g_yamlresume_url_str {://}
    { \url{#2} }
    { \url{http://#2} }
}
\ExplSyntaxOff

\name{Dr. Julia Schneider}{}
\title{Operations Manager | Prozessoptimierung \& Supply Chain}
\phone[mobile]{+49 89 12345678}
\email{julia.schneider@example.com}
\homepage{https://www.julia-schneider.de}

\address{Maximilianstraße 12, München, Bayern, Deutschland, 80333}{}{}

\extrainfo{{\small \faLinkedin}\ \href{https://www.linkedin.com/in/julia-schneider-ops}{@julia-schneider-ops} {} {} {} • {} {} {} 
{\small \faTwitter}\ \href{https://twitter.com/j\_schneider\_ops}{@j\_schneider\_ops}}

\begin{document}

\maketitle

\section{Basisinformationen}

\cvline{}{\begin{itemize}
\item Operations Manager mit über 10 Jahren Erfahrung in der Optimierung von Produktions- und Logistikprozessen
\item Nachweisliche Erfolge bei der Steigerung der Effizienz und Senkung der Betriebskosten durch Lean-Management und Six Sigma
\item Expertise in der Führung interdisziplinärer Teams und der Umsetzung strategischer Projekte
\item Fundierte Kenntnisse in SAP ERP, Supply Chain Management und kontinuierlicher Verbesserung (KVP)
\end{itemize}}
\section{Ausbildung}

\cventry{Okt. 2010–Sept. 2012}
        {Master, Wirtschaftsingenieurwesen, Punkte: 1,8}
        {Technische Universität München}
        {\url{https://www.tum.de}}
        {}
        {\begin{itemize}
\item Schwerpunkt auf der Optimierung von Wertschöpfungsketten und betrieblichen Prozessen
\item Praktische Erfahrung durch Projekte mit Industriepartnern, z. B. zur Lagerhaltungsoptimierung
\item Ausgezeichnet mit dem Preis der Fakultät für herausragende Masterarbeit
\end{itemize}
\textbf{Kurse}: Produktionsmanagement, Logistik und Supply Chain Management, Operations Research, Qualitätsmanagement, Projektmanagement, Controlling und Finanzen}
\section{Berufserfahrung}

\cventry{Jan. 2020–Heute}
        {Senior Operations Manager}
        {Nordlicht Logistik GmbH}
        {\url{https://www.nordlicht-logistik.de}}
        {}
        {\begin{itemize}
\item Leitung der Betriebsabteilung mit 45 Mitarbeitern und Verantwortung für ein Budget von 8 Mio. EUR
\item Einführung eines neuen Lagerverwaltungssystems (WMS) zur Steigerung der Lagerumschlagsgeschwindigkeit um 25 \%
\item Implementierung von Lean-Management-Prinzipien, wodurch die Durchlaufzeiten um 15 \% reduziert wurden
\item Aufbau eines KPI-Dashboards zur Echtzeitüberwachung der operativen Leistung
\end{itemize}
\textbf{Schlüsselwörter}: Lagerverwaltung, Lean Management, Budgetverantwortung, KPI-Reporting}

\cventry{März 2017–Dez. 2019}
        {Operations Manager}
        {TechFlow Solutions AG}
        {\url{https://www.techflow-solutions.de}}
        {}
        {\begin{itemize}
\item Verantwortung für die Auftragsabwicklung und Produktionsplanung in einem internationalen Umfeld
\item Optimierung der Lieferkette durch Einführung von S\&OP-Prozessen, Reduzierung der Lieferzeiten um 20 \%
\item Führung eines Teams von 20 Mitarbeitern und Einführung regelmäßiger Schulungen zu Qualitätsstandards
\item Zusammenarbeit mit Zulieferern zur Sicherstellung der Materialverfügbarkeit und Kostensenkung
\end{itemize}
\textbf{Schlüsselwörter}: Produktionsplanung, S\&OP, Lieferkettenoptimierung, Qualitätsmanagement}

\cventry{Juli 2014–Feb. 2017}
        {Operations Manager}
        {Rhein-Main Produktion GmbH}
        {\url{https://www.rhein-main-produktion.de}}
        {}
        {\begin{itemize}
\item Steuerung der täglichen Betriebsabläufe in einem Produktionswerk mit 120 Mitarbeitern
\item Einführung von Kanban-Systemen zur Bestandsreduzierung um 30 \%
\item Durchführung von Wertstromanalysen und Ableitung von Maßnahmen zur Effizienzsteigerung
\item Enge Zusammenarbeit mit der Geschäftsleitung zur Entwicklung langfristiger Betriebsstrategien
\end{itemize}
\textbf{Schlüsselwörter}: Kanban, Wertstromanalyse, Betriebssteuerung, Effizienzsteigerung}
\section{Sprachen}

\cvline{Deutsch}{Muttersprache oder zweisprachig \hfill \textbf{Schlüsselwörter}: Muttersprache}
\cvline{Englisch}{Verhandlungssichere Kenntnisse \hfill \textbf{Schlüsselwörter}: TOEFL iBT 110, Verhandlungsführung}
\section{Fähigkeiten}

\cvline{Prozessoptimierung}{Experte \hfill \textbf{Schlüsselwörter}: Lean Management, Six Sigma, Wertstromanalyse, KVP, Kanban}
\cvline{Supply Chain Management}{Erfahren \hfill \textbf{Schlüsselwörter}: S\&OP, Lieferkettenoptimierung, Lagerverwaltung, Beschaffung}
\cvline{Teamführung}{Erfahren \hfill \textbf{Schlüsselwörter}: Mitarbeiterentwicklung, Konfliktmanagement, Changemanagement, Coaching}
\cvline{SAP ERP}{Erfahren \hfill \textbf{Schlüsselwörter}: SAP MM, SAP PP, SAP WM}
\cvline{Projektmanagement}{Erfahren \hfill \textbf{Schlüsselwörter}: PMP, Agile Methoden, Stakeholder-Management, Risikomanagement}
\section{Auszeichnungen}

\cventry{Nov. 2022}
        {Nordlicht Logistik GmbH}
        {Best Operations Excellence Award}
        {}
        {}
        {Auszeichnung für die erfolgreiche Implementierung eines ganzheitlichen Lean-Management-Programms, das zu einer jährlichen Kosteneinsparung von 1,2 Mio. EUR führte.}

\cventry{Okt. 2012}
        {Technische Universität München}
        {Preis der Fakultät für herausragende Masterarbeit}
        {}
        {}
        {Für die Masterarbeit zum Thema "Optimierung von Lagerhaltungsstrategien mittels simulationsbasierter Analyse".}
\section{Zertifikate}

\cventry{Juni 2016}
        {TÜV Rheinland}
        {Lean Six Sigma Green Belt}
        {\url{https://www.tuv.com/lean-six-sigma}}
        {}
        {}

\cventry{Sept. 2018}
        {Project Management Institute}
        {Project Management Professional (PMP)}
        {\url{https://www.pmi.org/certifications/project-management-pmp}}
        {}
        {}

\cventry{März 2015}
        {SAP SE}
        {SAP Certified Application Associate - Supply Chain}
        {\url{https://www.sap.com/certification}}
        {}
        {}
\section{Projekte}

\cventry{März 2021–Nov. 2021}
        {Implementierung eines modernen WMS zur Optimierung der Lagerprozesse}
        {Einführung eines WMS}
        {\url{https://www.nordlicht-logistik.de/projekte/wms}}
        {}
        {\begin{itemize}
\item Auswahl und Einführung eines neuen Warehouse-Management-Systems bei Nordlicht Logistik
\item Analyse der bestehenden Prozesse und Definition von Anforderungen
\item Schulung der Mitarbeiter und Begleitung der Inbetriebnahme
\item Erzielung einer 25\%igen Steigerung der Lagerumschlagsgeschwindigkeit
\end{itemize}
\textbf{Schlüsselwörter}: WMS, Logistik, Prozessoptimierung}
\section{Interessen}

\cvline{Sport}{Laufen, Radfahren, Wandern}
\cvline{Logistik und Technologie}{Industrie 4.0, Automatisierung, Nachhaltigkeit}
\section{Ehrenamtliche Tätigkeiten}

\cventry{Jan. 2019–Heute}
        {Ehrenamtlicher Fachgruppenleiter}
        {Bundesverband Materialwirtschaft, Einkauf und Logistik e.V. (BME)}
        {\url{https://www.bme.de}}
        {}
        {\begin{itemize}
\item Leitung der regionalen Fachgruppe München zum Thema "Digitale Supply Chain"
\item Organisation von Fachvorträgen und Networking-Events für über 100 Mitglieder
\item Förderung des Wissensaustauschs zwischen Industrie und Wissenschaft
\end{itemize}}
    
\end{document}